\documentclass{article}

% The course guidelines require the NeurIPS LaTeX template.
% The provided Overleaf link currently points to the NeurIPS 2026 source.
\usepackage[preprint]{neurips_2026}

\usepackage[utf8]{inputenc}
\usepackage[T1]{fontenc}
\usepackage{hyperref}
\usepackage{url}
\usepackage{booktabs}
\usepackage{amsfonts}
\usepackage{nicefrac}
\usepackage{microtype}
\usepackage{xcolor}
\usepackage{graphicx}

\title{Predicting NBA Player Scoring Performance Using Historical Statistics}

\author{%
  Alec Choi \\
  NetID: alc409 \\
  Section ID: TBD \\
  \And
  Andre Sumilang \\
  NetID: aas504 \\
  Section ID: TBD \\
  \And
  Ryan Shentu \\
  NetID: rs2281 \\
  Section ID: TBD \\
}

\begin{document}

\maketitle

\begin{abstract}
This project studies whether NBA player scoring performance, measured as points per game, can be predicted from historical player-season statistics such as playing time, shot volume, shooting efficiency, free throw volume, and box score contributions. Using the NBA subset of the Kaggle NBA Stats (1947-present) dataset snapshot downloaded on April 30, 2026, we compare linear and tree-based regression models and evaluate their performance using MAE, MSE, RMSE, and $R^2$. Code: \url{https://github.com/alecchoi8/Intro-to-DS---Final-Project} Video: \url{TODO}
\end{abstract}

\section{Introduction}

Sports analytics has become central to evaluating player performance and understanding team offense. This project asks whether a player's points per game can be predicted from player-season statistics available in historical NBA data.

\section{Related Work}

TODO: Discuss basketball analytics, player performance modeling, Basketball-Reference style box score data, and prior use of regression models for sports prediction.

\section{Data}

TODO: Describe the Kaggle NBA Stats (1947-present) dataset, the April 30, 2026 snapshot, NBA-only filtering, player-season observation level, and manual aggregation of multi-team player seasons.

\section{Method}

TODO: Describe preprocessing, feature selection, train/test split, and model pipelines. Required models include Linear Regression, Decision Tree Regression, and Random Forest Regression. Optional comparison models may include regularized regression or gradient boosting.

\section{Evaluation}

TODO: Report MAE, MSE, RMSE, and $R^2$ for each model. Include a performance table and actual-versus-predicted plot.

\section{Discussion}

TODO: Interpret model results, key predictors, feature importances, and what the results suggest about NBA scoring output.

\section{Limitations}

TODO: Discuss era effects, missing historical statistics, rule changes, incomplete modern-season data if relevant, and the fact that some scoring component features are closely related to the target.

\section{Conclusion}

TODO: Summarize findings and future extensions.

\bibliographystyle{plain}
\bibliography{references}

\end{document}
